\section{Introduction}

Every frontier-based navigation agent faces the same dilemma: it must decide where to explore next using a scoring function that knows nothing about whether its previous decisions were any good. A frontier evaluated at step one is assessed with the same model, the same confidence, and the same exploration-exploitation balance as one evaluated at step two hundred, regardless of how many wrong rooms the agent has already walked into. The search strategy is \emph{static}, and in Embodied Question Answering (EQA), where the agent must navigate an unfamiliar 3D environment to gather evidence for answering a natural language question, this rigidity is especially costly. The agent must jointly solve \emph{what} to search for, \emph{where} to look, and \emph{how long} to keep exploring, yet it has no mechanism to learn from its own mistakes within an episode.

Recent memory-augmented methods~\cite{huang20243dmem,yang2024ramem,wang2026memoryexplorer} have made impressive progress on the reasoning side of EQA, showing that richer memory representations lead to better answers. But these methods inherit whatever observations navigation provides. If the agent wastes steps exploring irrelevant rooms, the memory is impoverished no matter how sophisticated the retrieval. \textcolor{red}{A further constraint comes from how the task itself is posed. Long-horizon exploration benchmarks conventionally treat every object named in an instruction as a terminal goal, to be reached in turn. We keep the data and read the instruction differently, as \emph{landmark-conditioned} navigation: an instruction specifies one target together with the landmarks describing the route toward it, and only the target must be reached. The landmarks are instrumental rather than terminal, and their value is informational. Each is a prediction, supplied by the instruction, of what the agent should encounter on the way. Verifying them requires no additional annotation, since this relational context is already present in the description goals, and it recovers a signal that multi-goal formulations discard.}

We introduce a simple mechanism that breaks this limitation: \emph{predict before you explore, then learn from your mistakes.} Before moving through a frontier, the agent uses a vision-language model to predict the semantic layout beyond it\textcolor{red}{, conditioned on the landmark the route expects at that stage}: what room type and object categories are likely to appear. After arriving, it compares this prediction against what it actually observes. The discrepancy, measured as KL divergence, produces a \emph{surprise signal} that reveals how well the agent understands the current environment. High surprise means the agent's model is miscalibrated; the rational response is to broaden exploration. Low surprise means the model has converged; the rational response is to narrow the search toward the goal. This exploration-exploitation shift emerges automatically from the surprise dynamics, with no explicit scheduling and no environment-specific training.

This predict-then-correct loop is made practical by the capabilities of modern vision-language models, which can produce meaningful semantic layout predictions from a single egocentric frontier view. Even when these predictions are imperfect (and they frequently are), the \emph{pattern of errors} is informative: systematic misprediction in one environment tells the agent something that no individual prediction can.

We instantiate this idea in \textbf{PEC-Nav} (Predict-Explore-Correct Navigation), a language-driven EQA pipeline whose search module forms a closed feedback loop. Given a natural language instruction, the pipeline extracts a structured navigation goal \textcolor{red}{and an ordered landmark sequence }via LLM parsing, then applies the PEC loop: a semantic predictor anticipates the layout beyond each frontier, an uncertainty-aware policy selects frontiers by balancing goal relevance against prediction confidence, and a correction module updates both the predictor's confidence scaling and the policy's exploration-exploitation weights online (Fig.~\ref{fig:architecture}). Observations gathered during navigation populate an episodic memory from which the agent answers follow-up questions. Because the memory and QA modules are held constant across all experiments, any change in answer quality is directly attributable to search quality\textcolor{red}{, isolating search as the sole experimental variable}.

We evaluate PEC-Nav on \textcolor{red}{the description goals of }GOAT-Bench Val Unseen~\cite{khanna2024goatbench} (36 scenes, 345 episodes). \textcolor{red}{Because the reformulation changes which objects must be reached, published results on the multi-goal task are reported for context rather than as head-to-head comparisons.} Using a 4B-parameter VLM on a single consumer GPU, PEC-Nav achieves 29.74\% SR and 20.88 SPL, competitive with methods that rely on models orders of magnitude larger. The correction ablation provides the sharpest evidence for our thesis: disabling the feedback loop while keeping the predictor active \emph{degrades} performance to 12.75\% SR, worse than not predicting at all. Re-enabling correction yields a 133\% relative improvement, confirming that it is prediction \emph{errors}, not predictions, that drive effective search.

Our contributions are:
\begin{itemize}
    \item \textcolor{red}{\textbf{Landmark-conditioned navigation,} a reformulation in which route landmarks are instrumental predictions to be verified rather than terminal goals to be reached, recovering a supervisory signal from existing data at no annotation cost.}
    \item \textbf{The Predict-Explore-Correct loop,} the first self-calibrating frontier search mechanism for embodied navigation. The agent predicts semantic layouts beyond frontiers, measures surprise upon arrival, and uses accumulated surprise to adapt both predictor confidence and exploration-exploitation balance within a single episode.
    \item \textbf{A diagnostic EQA pipeline} in which memory and reasoning are held fixed, isolating search as the sole experimental variable.
    \item \textbf{Empirical validation} showing that self-calibrating search outperforms static alternatives in navigation efficiency and downstream QA, with within-episode calibration dynamics measured directly from correction logs.
\end{itemize}

